\chapter{การเชื่อมต่อฮาร์ดแวร์และโปรโตคอล Modbus RTU}
\label{chap:ch02}

\section*{แผนบริหารการสอนประจำบทที่ 2}
\addcontentsline{toc}{section}{แผนบริหารการสอนประจำบทที่ 2}

\begin{tcolorbox}[
    enhanced, breakable,
    colback=white, colframe=rbruNavy,
    fonttitle=\bfseries\color{white},
    title={1. หัวข้อเนื้อหาประจำบท},
    arc=1.5mm, boxrule=0.8pt, leftrule=3.5mm
]
\begin{enumerate}
    \item ฟิสิกส์ของการส่งสัญญาณเชิงผลต่างและคุณลักษณะของสายส่ง RS485
    \item สถาปัตยกรรมหัววัด Soil Parameter Tester และการจ่ายพลังงานผ่าน USB-C OTG
    \item โครงสร้างเฟรมและการกำหนดจังหวะเวลาของโปรโตคอล Modbus RTU
    \item ทฤษฎีการตรวจสอบความสมบูรณ์ของข้อมูลด้วยผลรวมวงรอบ 16 บิต (CRC-16 Checksum)
    \item การเข้าสายสัญญาณและผังรีจิสเตอร์ข้อมูลตรวจวัด pH อุณหภูมิ และความชื้น
\end{enumerate}
\end{tcolorbox}

\begin{tcolorbox}[
    enhanced, breakable,
    colback=white, colframe=rbruNavy,
    fonttitle=\bfseries\color{white},
    title={2. วัตถุประสงค์เชิงพฤติกรรม},
    arc=1.5mm, boxrule=0.8pt, leftrule=3.5mm
]
เมื่อศึกษาบทเรียนนี้จบแล้ว ผู้เรียนสามารถ
\begin{enumerate}
    \item อธิบายหลักการหักล้างสัญญาณรบกวนร่วม (Common-Mode Noise Rejection) ของสายส่ง RS485 ได้อย่างถูกต้อง
    \item ต่อวงจรสายสัญญาณระหว่างหัววัดดินกับสายแปลง USB-RS485 Type-C ได้อย่างแม่นยำ
    \item ถอดรหัสโครงสร้างเฟรมคำสั่งและเฟรมข้อมูลตอบกลับของ Modbus RTU ได้อย่างสมบูรณ์
    \item เขียนและอธิบายอัลกอริทึมการคำนวณรหัสตรวจสอบความถูกต้อง CRC-16 ได้อย่างถูกต้อง
\end{enumerate}
\end{tcolorbox}

\begin{tcolorbox}[
    enhanced, breakable,
    colback=white, colframe=rbruNavy,
    fonttitle=\bfseries\color{white},
    title={3. วิธีสอนและกิจกรรมการเรียนการสอน},
    arc=1.5mm, boxrule=0.8pt, leftrule=3.5mm
]
\begin{enumerate}
    \item การบรรยายเชิงทฤษฎีควบคู่กับการใช้ออสซิลโลสโคปแสดงรูปคลื่นสัญญาณเชิงผลต่าง
    \item การสาธิตการส่งเฟรมคำสั่ง Modbus ผ่านโปรแกรม Serial Monitor
    \item กิจกรรมฝึกปฏิบัติการต่อสายฮาร์ดแวร์จริงและทดสอบการสื่อสารบนสมาร์ตโฟน
\end{enumerate}
\end{tcolorbox}

\begin{tcolorbox}[
    enhanced, breakable,
    colback=white, colframe=rbruNavy,
    fonttitle=\bfseries\color{white},
    title={4. สื่อการเรียนการสอน},
    arc=1.5mm, boxrule=0.8pt, leftrule=3.5mm
]
\begin{enumerate}
    \item เอกสารคำสอนวิชาการเชื่อมต่อระบบสมองกลฝังตัวทางการเกษตร
    \item ชุดหัววัด Soil Parameter Tester (SoilpH-HT) และสายแปลง USB Type-C OTG
    \item เครื่องออสซิลโลสโคปดิจิทัล 2 ช่องสัญญาณสำหรับการตรวจวัดรูปคลื่น RS485
\end{enumerate}
\end{tcolorbox}

\begin{tcolorbox}[
    enhanced, breakable,
    colback=white, colframe=rbruNavy,
    fonttitle=\bfseries\color{white},
    title={5. การวัดและประเมินผล},
    arc=1.5mm, boxrule=0.8pt, leftrule=3.5mm
]
\begin{enumerate}
    \item การประเมินผลการคำนวณบิต CRC-16 ในแบบฝึกหัด
    \item การประเมินทักษะการต่อสายสัญญาณและการแก้ปัญหาข้อผิดพลาดการเชื่อมต่อ
    \item การสังเกตความละเอียดรอบคอบและมาตรฐานความปลอดภัยทางไฟฟ้า
\end{enumerate}
\end{tcolorbox}

\clearpage

% ==============================================================================
% ภาพประกอบเกริ่นนำประจำบทที่ 2 (Introductory Infographic)
% ==============================================================================
\begin{figure}[H]
\centering
\begin{tikzpicture}[
    box/.style={rectangle, rounded corners=2mm, draw=rbruNavy, line width=1.0pt, fill=white, drop shadow={shadow xshift=1pt, shadow yshift=-1.5pt, opacity=0.15}, align=center, inner sep=6pt},
    header/.style={rectangle, rounded corners=1.5mm, draw=rbruNavy, fill=rbruNavy, text=white, font=\bfseries\small, align=center, inner sep=4pt},
    wire/.style={line width=1.4pt},
    arrow/.style={-{Stealth[scale=1.1]}, line width=1.2pt, draw=rbruBlue}
]
    % โหนดที่ 1: สมาร์ตโฟน
    \node[box, fill=cyan!10, minimum width=3.2cm] (phone) {
        \textbf{สมาร์ตโฟนแอนดรอยด์}\\[2pt]
        \footnotesize
        $\bullet$ USB Host Mode\\
        $\bullet$ USB OTG Power\\
        $\bullet$ SoilpH-HT App
    };
    \node[header, above=2pt of phone] {1. ฝั่งควบคุมหลัก};

    % โหนดที่ 2: ชิปแปลงสัญญาณ USB-RS485
    \node[box, fill=rbruGold!10, minimum width=3.4cm, right=1.0cm of phone] (adapter) {
        \textbf{สายแปลง USB-RS485}\\[2pt]
        \footnotesize
        $\bullet$ CH340 / FTDI Bridge\\
        $\bullet$ MAX485 Transceiver\\
        $\bullet$ DE/RE Auto Direction
    };
    \node[header, fill=rbruGold!90!black, above=2pt of adapter] {2. การแปลงสัญญาณ};

    % โหนดที่ 3: สายส่งเชิงผลต่าง
    \node[box, fill=rbruNavy!10, minimum width=3.2cm, right=1.0cm of adapter] (bus) {
        \textbf{สายส่งคู่ตีเกลียว}\\[2pt]
        \footnotesize
        $\bullet$ สาย A+ (Non-inverting)\\
        $\bullet$ สาย B- (Inverting)\\
        $\bullet$ $Z_0 = 120\,\Omega$
    };
    \node[header, fill=rbruNavy, above=2pt of bus] {3. สายส่งเชิงผลต่าง};

    % โหนดที่ 4: หัววัดดินจริง
    \node[box, fill=green!10, minimum width=3.2cm, right=1.0cm of bus] (probe) {
        \textbf{Soil Parameter Tester}\\[2pt]
        \footnotesize
        $\bullet$ Modbus RTU Slave\\
        $\bullet$ วัด pH, อุณหภูมิ, ความชื้น\\
        $\bullet$ ขั้วเข็มสแตนเลส 316
    };
    \node[header, fill=green!60!black, above=2pt of probe] {4. เซนเซอร์ภาคสนาม};

    % สายเชื่อมต่อ
    \draw[arrow] (phone) -- node[above, font=\scriptsize\bfseries] {USB-C} (adapter);
    \draw[arrow] (adapter) -- node[above, font=\scriptsize\bfseries] {TTL/UART} (bus);
    \draw[arrow] (bus) -- node[above, font=\scriptsize\bfseries] {RS485} (probe);

    % กรอบสรุปด้านล่าง
    \node[rectangle, rounded corners=2mm, draw=rbruGold!80, fill=rbruGold!8, line width=0.8pt, minimum width=14.5cm, minimum height=0.85cm, below=0.7cm of adapter, xshift=1.5cm, align=center] (proto) {
        \small\textbf{โปรโตคอล Modbus RTU} อัตราเร็ว 4,800 bps, 8 Data Bits, No Parity, 1 Stop Bit, ตรวจสอบด้วย CRC-16
    };
    \draw[arrow, dashed] (adapter.south) -- ++(0,-0.4) -| (proto.north);
\end{tikzpicture}
\caption{ผังการเชื่อมโยงทางกายภาพและการไหลของสัญญาณดิจิทัลผ่านพอร์ต USB Type-C OTG สู่หัววัด Soil Parameter Tester}
\label{fig:ch02_hardware_interconnect}
\end{figure}

\section{ฟิสิกส์ของการส่งสัญญาณเชิงผลต่าง RS485}

การส่งสัญญาณดิจิทัลในแปลงเกษตรกรรมต้องเผชิญกับสัญญาณรบกวนแม่เหล็กไฟฟ้า (Electromagnetic Interference หรือ EMI) จากมอเตอร์ปั๊มน้ำ เครื่องยนต์กลเกษตร และฟ้าแลบ มาตรฐานสายส่งเชิงผลต่าง \textbf{RS485 (TIA/EIA-485-A)} จึงถูกเลือกใช้ด้วยคุณสมบัติการหักล้างสัญญาณรบกวนร่วมที่ยอดเยี่ยม

\subsection{หลักการส่งสัญญาณเชิงผลต่าง (Differential Signaling)}
สัญญาณดิจิทัลบนสายส่ง RS485 ไม่ได้วัดเทียบกับกราวด์ แต่วัดจากความต่างศักย์ระหว่างสายนำสัญญาณสองเส้น ได้แก่ เส้น A (Non-inverting) และเส้น B (Inverting) ดังสมการ

\begin{equation}
V_{AB} = V_A - V_B
\end{equation}

ตามมาตรฐาน สถานะลอจิกถูกนิยามดังนี้
\begin{itemize}
    \item \textbf{ลอจิก 1 (Mark / Off)} เมื่อ $V_A - V_B < -200 \text{ mV}$ (โดยทั่วไป $V_B > V_A$)
    \item \textbf{ลอจิก 0 (Space / On)} เมื่อ $V_A - V_B > +200 \text{ mV}$ (โดยทั่วไป $V_A > V_B$)
\end{itemize}

เมื่อมีสัญญาณรบกวนจากภายนอกเหนี่ยวนำเข้ามาในสายคู่เกลียว (Twisted Pair) สัญญาณรบกวนจะปรากฏบนสายทั้งสองเส้นในปริมาณที่เท่ากัน ($V_{\text{noise}}$) วงจรรับสัญญาณเชิงผลต่างจะทำการลบสัญญาณออกจากกัน ส่งผลให้สัญญาณรบกวนถูกหักล้างหมดไป ดังนี้

\begin{equation}
(V_A + V_{\text{noise}}) - (V_B + V_{\text{noise}}) = V_A - V_B = V_{AB}
\end{equation}

\section{สถาปัตยกรรมหัววัดและผังรีจิสเตอร์ข้อมูล Modbus}

หัววัด Soil Parameter Tester (SoilpH-HT) ผลิตจากสแตนเลสสตีลเกรด 316L ทนทานต่อการกัดกร่อนจากกรดและเกลือในดิน การเชื่อมต่อผ่านสายแปลง USB-RS485 Type-C ช่วยให้สมาร์ตโฟนสามารถจ่ายไฟเลี้ยง +5V DC จากพอร์ต USB ไปเลี้ยงหัววัดและวงจรประมวลผลได้โดยตรง

\begin{table}[H]
\centering
\caption{การกำหนดรหัสสีสายสัญญาณและหน้าที่ของหัววัด Soil Parameter Tester}
\label{tab:wire_pinout}
\begin{tabularx}{\textwidth}{p{2.5cm}p{3.5cm}Y}
    \toprule
    \textbf{สีของสาย} & \textbf{ชื่อสัญญาณ} & \textbf{หน้าที่และการต่อใช้งาน} \\
    \midrule
    สีน้ำตาลหรือสีแดง & VCC (+5V ถึง +12V DC) & ต่อเข้ากับขั้วบวกแหล่งจ่ายไฟจากสายแปลง USB OTG \\
    สีดำ & GND (0V) & ต่อเข้ากับขั้วลบหรือกราวด์ร่วมของระบบ \\
    สีเหลือง & RS485 Data A+ & ขั้วสัญญาณข้อมูลแบบไม่กลับเฟส (Non-inverting) \\
    สีน้ำเงินหรือสีเขียว & RS485 Data B- & ขั้วสัญญาณข้อมูลแบบกลับเฟส (Inverting) \\
    \bottomrule
\end{tabularx}
\end{table}

\begin{table}[H]
\centering
\caption{ผังรีจิสเตอร์ข้อมูล Modbus Holding Registers สำหรับการตรวจวัดค่า pH อุณหภูมิ และความชื้น}
\label{tab:modbus_ph_registers}
\begin{tabularx}{\textwidth}{p{2.5cm}p{2.2cm}p{2.5cm}p{2.2cm}Y}
    \toprule
    \textbf{ที่อยู่รีจิสเตอร์} & \textbf{พารามิเตอร์} & \textbf{ประเภทข้อมูล} & \textbf{ตัวคูณสเกล} & \textbf{หน่วยวัดฟิสิกส์} \\
    \midrule
    \texttt{0x0000} & ความชื้นดิน ($M_{\text{soil}}$) & 16-bit Unsigned & หารด้วย 10.0 & ร้อยละเชิงปริมาตร (\% VWC) \\
    \texttt{0x0001} & อุณหภูมิดิน ($T_{\text{soil}}$) & 16-bit Signed & หารด้วย 10.0 & องศาเซลเซียส ($^\circ\text{C}$) \\
    \texttt{0x0002} & สภาพนำไฟฟ้า ($EC_{\text{raw}}$) & 16-bit Unsigned & จำนวนเต็ม & ไมโครซีเมนต์ต่อเซนติเมตร ($\mu\text{S/cm}$) \\
    \texttt{0x0003} & ความเป็นกรด-ด่าง ($\text{pH}_{\text{raw}}$) & 16-bit Unsigned & หารด้วย 10.0 & หน่วย pH ($0.0 - 14.0$) \\
    \bottomrule
\end{tabularx}
\end{table}

\section{โครงสร้างเฟรมและการกำหนดจังหวะเวลาของโปรโตคอล Modbus RTU}

การสื่อสารระหว่างสมาร์ตโฟนกับหัววัดดินทำงานที่อัตราความเร็ว 4,800 บิตต่อวินาที (Baud Rate 4800, 8 Data Bits, No Parity, 1 Stop Bit) โดยมีช่วงเวลาของ 1 ตัวอักษร ($t_{\text{char}}$) เท่ากับ

\begin{equation}
t_{\text{char}} = \frac{10 \text{ บิต}}{4,800 \text{ บิต/วินาที}} \approx 2.083 \text{ มิลลิวินาที}
\end{equation}

ตามมาตรฐาน Modbus RTU การเริ่มต้นและสิ้นสุดเฟรมข้อมูลจะต้องมีช่วงเวลาเงียบว่าง (Inter-frame Silent Interval) ไม่น้อยกว่า 3.5 ตัวอักษร ($t_{3.5} \ge 7.29 \text{ มิลลิวินาที}$) เพื่อให้ตัวรับสัญญาณสามารถตรวจจับจุดเริ่มต้นของเฟรมใหม่ได้อย่างถูกต้อง

\begin{figure}[H]
\centering
\begin{tikzpicture}[
    node distance=1.5cm,
    byte/.style={rectangle, draw=rbruNavy, fill=rbruCodeBg, thick, minimum width=1.5cm, minimum height=0.8cm, align=center, font=\small},
    timing/.style={dashed, draw=gray, thick}
]
    \node[byte, fill=orange!20] (silent1) {Silent\\$\ge 3.5t$};
    \node[byte, fill=cyan!20, right=0pt of silent1] (addr) {Slave ID\\(1 Byte)};
    \node[byte, fill=cyan!30, right=0pt of addr] (func) {Function\\(1 Byte)};
    \node[byte, fill=cyan!40, right=0pt of func] (data) {Data Bytes\\(N Bytes)};
    \node[byte, fill=rbruGold!30, right=0pt of data] (crc) {CRC-16\\(2 Bytes)};
    \node[byte, fill=orange!20, right=0pt of crc] (silent2) {Silent\\$\ge 3.5t$};

    \draw[timing] ($(silent1.north west) + (0, 0.4)$) -- ($(silent1.south west) + (0, -0.4)$);
    \draw[timing] ($(silent2.north east) + (0, 0.4)$) -- ($(silent2.south east) + (0, -0.4)$);
\end{tikzpicture}
\caption{โครงสร้างจังหวะเวลาของเฟรมข้อมูลในโปรโตคอล Modbus RTU}
\label{fig:modbus_timing}
\end{figure}

\section{ทฤษฎีการตรวจสอบความสมบูรณ์ด้วยผลรวมวงรอบ 16 บิต (CRC-16)}

เพื่อป้องกันความผิดพลาดจากการรบกวนสัญญาณในสายส่ง โปรโตคอล Modbus RTU กำหนดให้ใช้รหัสตรวจสอบความถูกต้องแบบผลรวมวงรอบ 16 บิต (Cyclic Redundancy Check หรือ CRC-16) โดยใช้พหุนามกำเนิดมาตรฐาน

\begin{equation}
G(x) = x^{16} + x^{15} + x^2 + 1 \quad \implies \quad \mathtt{0xA001} \text{ (Bit-reversed)}
\end{equation}

\subsection{ขั้นตอนวิธีคำนวณ CRC-16 เชิงตัวเลข}
\begin{enumerate}
    \item กำหนดค่าเริ่มต้นของรีจิสเตอร์ CRC ขนาด 16 บิต เท่ากับ $\mathtt{0xFFFF}$
    \item นำไบต์ข้อมูลไบต์แรกมาทำการ Exclusive-OR (XOR) กับไบต์ต่ำของรีจิสเตอร์ CRC
    \item เลื่อนบิตรีจิสเตอร์ CRC ไปทางขวา 1 บิต (Shift Right)
    \item ตรวจสอบบิตที่หลุดออกมา (Least Significant Bit หรือ LSB)
    หากบิตที่หลุดออกมาเป็น 1 ให้ทำการ XOR รีจิสเตอร์ CRC กับพหุนาม $\mathtt{0xA001}$
    หากบิตที่หลุดออกมาเป็น 0 ไม่ต้องทำการ XOR
    \item ทำซ้ำขั้นตอนที่ 3 และ 4 จนครบ 8 ครั้ง (ครบ 1 ไบต์ข้อมูล)
    \item ทำซ้ำขั้นตอนที่ 2 ถึง 5 สำหรับไบต์ข้อมูลถัดไปจนครบทุกไบต์ในเฟรม
    \item ผลลัพธ์สุดท้ายคือค่า CRC-16 โดยต้องส่งไบต์ต่ำออกไปก่อนไบต์สูงตามมาตรฐาน Little-Endian
\end{enumerate}

\section{สถาปัตยกรรมการเชื่อมต่อ USB-C OTG สำหรับเครื่องวัดพารามิเตอร์ดิน}

ในการประยุกต์ใช้งานภาคสนามร่วมกับสมาร์ตโฟนแอนดรอยด์ หัววัดดินอุตสาหกรรมจำเป็นต้องเชื่อมต่อผ่านสายแปลง USB-RS485 Type-C ที่รองรับโหมด USB On-The-Go (USB OTG) โดยมีองค์ประกอบการทำงานที่สำคัญดังนี้

\subsection{ผังการเข้าสายสัญญาณและการจ่ายพลังงาน}
หัววัดดินอุตสาหกรรมใช้สายเชื่อมต่อมาตรฐาน 4 เส้น ซึ่งต้องต่อเข้ากับขั้วของโมดูลแปลงสัญญาณ USB-RS485 ดังนี้
\begin{enumerate}
    \item \textbf{สายสีแดง (VCC)} จ่ายแรงดันไฟฟ้ากระแสตรง $5 - 30\text{ โวลต์}$ โดยสามารถดึงไฟเลี้ยง $5\text{ โวลต์}$ จากพอร์ต USB Type-C ของสมาร์ตโฟนผ่านวงจร Step-Up Converter ได้โดยตรง
    \item \textbf{สายสีดำ (GND)} กราวด์ร่วมของระบบไฟฟ้า
    \item \textbf{สายสีเหลือง (RS485 A+)} สัญญาณผลต่างขั้วบวก (Non-inverting Data Line)
    \item \textbf{สายสีน้ำเงินหรือสีเขียว (RS485 B-)} สัญญาณผลต่างขั้วลบ (Inverting Data Line)
\end{enumerate}

\subsection{การกำหนดค่าสิทธิ์ในระบบแอนดรอยด์ (Android USB Host Permissions)}
ระบบปฏิบัติการแอนดรอยด์กำหนดให้แอปพลิเคชันต้องขอสิทธิ์การเข้าถึงอุปกรณ์ USB จากผู้ใช้ โดยในไฟล์ \texttt{AndroidManifest.xml} ต้องประกาศการใช้งานฟีเจอร์ \texttt{android.hardware.usb.host} และผูกเข้ากับไฟล์ตัวกรอง \texttt{device\_filter.xml} ซึ่งระบุหมายเลขประจำตัวผู้ผลิต (Vendor ID หรือ VID) ของชิปแปลง USB-Serial สากล ได้แก่
\begin{itemize}
    \item ชิป WCH CH340 / CH341 (Vendor ID \texttt{0x1A86})
    \item ชิป Silicon Labs CP2102 / CP2104 (Vendor ID \texttt{0x10C4})
    \item ชิป FTDI FT232R (Vendor ID \texttt{0x0403})
    \item ชิป Prolific PL2303 (Vendor ID \texttt{0x067B})
\end{itemize}

\subsection{ขั้นตอนวิธีตรวจจับอัตราเร็วอัตโนมัติ (Auto-Baud Rate Detection)}
หัววัดดินในท้องตลาดมักตั้งค่าอัตราเร็วเริ่มต้นมาที่ $4,800\text{ bps}$ หรือ $9,600\text{ bps}$ ระบบแอปพลิเคชัน \texttt{SoilpH-HT} จึงได้รับการออกแบบให้อัลกอริทึมทดลองส่งเฟรมคำถามไปยังหัววัดด้วยอัตราเร็วแรก หากไม่ได้รับการตอบกลับหรือค่าผลรวมตรวจสอบ CRC-16 ผิดพลาดติดต่อกันเกิน 3 ครั้ง ระบบจะสลับอัตราเร็วไปยัง $9,600\text{ bps}$ โดยอัตโนมัติ ช่วยลดความยุ่งยากของผู้ใช้งานที่ไม่มีประสบการณ์ด้านฮาร์ดแวร์

\subsection{ระบบเฝ้าระวังการสื่อสารและเชื่อมต่อซ้ำอัตโนมัติ (Watchdog \& Hotplug)}
เพื่อความทนทานต่อการใช้งานจริงในแปลงเกษตร ระบบประกอบด้วยวงจรเฝ้าระวังทางซอฟต์แวร์ (Communication Watchdog Timer) ขนาด 3 วินาที หากไม่มีข้อมูลส่งกลับมาจากหัววัดเกินกว่าช่วงเวลาดังกล่าว ระบบจะตัดการทำงานของพอร์ตชั่วคราวและสั่งเชื่อมต่อใหม่โดยอัตโนมัติ พร้อมทั้งมีตัวรับการแจ้งเตือนระดับระบบปฏิบัติการ (\texttt{BroadcastReceiver}) ดักจับเหตุการณ์ \texttt{ACTION\_USB\_ATTACHED} ทำให้เมื่อผู้ใช้เสียบสายโพรบเข้ากับพอร์ต Type-C แอปพลิเคชันจะตรวจพบและเริ่มอ่านค่าได้ทันทีโดยไม่ต้องเปิดปิดแอปพลิเคชันใหม่

\section{สรุปสาระสำคัญประจำบทที่ 2}

การเชื่อมต่อสมาร์ตโฟนเข้ากับเซนเซอร์ฮาร์ดแวร์จริงภาคสนาม จำเป็นต้องเข้าใจทั้งชั้นกายภาพ (Physical Layer) และชั้นโปรโตคอล (Data Link Layer) โดยสรุปสาระสำคัญได้ดังนี้
\begin{enumerate}
    \item \textbf{สายส่งเชิงผลต่าง RS485} ใช้สายสัญญาณคู่ตีเกลียววัดผลต่างศักย์ $V_{AB} = V_A - V_B$ ช่วยหักล้างสัญญาณรบกวนร่วม (Common-Mode Rejection) จากอุปกรณ์ไฟฟ้าในแปลงได้อย่างมีประสิทธิภาพ
    \item \textbf{โปรโตคอล Modbus RTU} สื่อสารแบบ Master-Slave ผ่านอัตราเร็ว 4,800 หรือ 9,600 bps โดยมีช่วงเวลาเงียบว่าง $\ge 3.5$ ตัวอักษร ($t_{3.5}$) เป็นตัวแบ่งเฟรมข้อมูล
    \item \textbf{การตรวจสอบความถูกต้องด้วย CRC-16} ทุกเฟรมข้อมูลต้องผ่านการคำนวณและตรวจสอบรหัสตรวจสอบข้อผิดพลาดขนาด 16 บิต เพื่อป้องกันข้อมูลเพี้ยนจากสัญญาณรบกวนก่อนส่งต่อไปยังชั้นประมวลผล
    \item \textbf{การควบคุมทิศทาง Half-Duplex (RTS)} ในการรับส่งข้อมูลแบบสลับทิศทาง จะต้องปลดสาย RTS ให้เป็น Low เพื่อให้ตัวรับสัญญาณ (Receiver) เปิดรับข้อมูลจากหัววัดเข้าสู่สมาร์ตโฟน
\end{enumerate}

ในบทถัดไป (บทที่ 3) จะนำข้อมูลดิบที่รับมาจากหัววัดผ่านโปรโตคอล Modbus RTU ไปประมวลผลและชดเชยค่าความคลาดเคลื่อนด้วยโครงข่ายประสาทเทียมอิงฟิสิกส์ (PINN)

\section*{คำถามทบทวนท้ายบทที่ 2}
\addcontentsline{toc}{section}{คำถามทบทวนท้ายบทที่ 2}

\begin{enumerate}
    \item จงอธิบายหลักการทำงานของสัญญาณเชิงผลต่าง (Differential Signaling) ในมาตรฐาน RS485 และเหตุผลที่สามารถหักล้างสัญญาณรบกวนร่วมได้
    \item เพราะเหตุใดในโปรโตคอล Modbus RTU จึงจำเป็นต้องกำหนดช่วงเวลาเงียบว่างระหว่างเฟรมไม่น้อยกว่า 3.5 ตัวอักษร ($t_{3.5}$)
    \item จงแสดงการคำนวณหาค่า CRC-16 ของชุดข้อมูล 2 ไบต์ ได้แก่ $[\mathtt{0x01}, \mathtt{0x03}]$ ตามขั้นตอนวิธีที่ระบุในบทเรียน
    \item หากเชื่อมต่อหัววัดดินเข้ากับสมาร์ตโฟนแล้วพบว่าไฟแสดงสถานะบนสายแปลงกะพริบแต่แอปพลิเคชันแจ้งเตือน CRC Error จงวิเคราะห์สาเหตุและแนวทางแก้ไข
\end{enumerate}
